% `tex.web` §298 names a character token with a category-code word and then
% the character:
%
%     letter: print("the letter ");
%     other_char: print("the character ");
%     …
%     othercases print_char(chr_cmd)  { and then `print(chr)` }
%
% §59's `print` sends anything outside the printable range through §48's `^^`
% form — the same one `print_ASCII` uses.  With `\catcode`\^=7`, `^^^A` is the
% character 30, which TeX names `^^^`; we pushed the control character itself,
% which is invisible in a terminal and absent from a log.
%
% trip.tex reaches this at command 824, where a subscript character is `^^A`.
\catcode`\{=1 \catcode`\}=2 \catcode`\$=3
\catcode`\^=7 \catcode1=8 \catcode2=7 \catcode200=11
\font\r=cmr10 \r \hsize100pt \parindent0pt \hbadness10000 \tolerance10000
\tracingcommands2 \tracingonline0
\textfont0=\r \scriptfont0=\r \scriptscriptfont0=\r
\textfont1=\r \scriptfont1=\r \scriptscriptfont1=\r
\textfont2=\r \scriptfont2=\r \scriptscriptfont2=\r
\textfont3=\r \scriptfont3=\r \scriptscriptfont3=\r
\fontdimen22\r=1pt
$x^^^Ay$
$x^^A y^^B z^^c8$
\tracingcommands0
\par\end
